\documentclass[11pt,a4paper]{article}
\usepackage[margin=0.85in,top=1in,bottom=1in]{geometry}
\usepackage{amsmath,amssymb,amsfonts}
\usepackage{graphicx}
\usepackage{float}
\usepackage{enumitem}
\usepackage{microtype}
\usepackage{caption}
\usepackage[hidelinks]{hyperref}
\usepackage[most,breakable]{tcolorbox}
\tcbuselibrary{breakable,skins}
\usepackage[utf8]{inputenc}
\usepackage[T1]{fontenc}
\usepackage{lmodern}
\captionsetup{font=small,labelfont=bf,skip=4pt}
\setlist{leftmargin=1.5em,topsep=2pt}

%% Breakable card — prevents content cutoff across pages
\newtcolorbox{solcard}[3]{
  breakable, enhanced,
  colback=white, colframe=gray!50, arc=2pt,
  title={\textbf{#1}\hfill\textsf{\small #3\,pts}},
  fonttitle=\normalsize\sffamily\bfseries,
  before upper={\small\textit{#2}\par\smallskip\hrule height 0.4pt\smallskip},
  top=4pt, bottom=6pt, left=7pt, right=7pt,
  parbox=false,
}

\title{\textbf{Thomson jumping ring --- Solution}}
\author{\normalsize Y23 (2023) — English translation}
\date{}
\begin{document}\maketitle\vspace{-0.5em}
\begin{solcard}{A1}{Write down an equation relating the current in the ring $I(t)$ and the magnetic field in the core $B_z(t)$.}{0.30}

Let's write Faraday's law of electromagnetic induction: $$\mathcal{E}_{ind}=-\cfrac{d\Phi}{dt} $$For choose direction current in the ring (counterclockwise viewed from above) the induced EMF equals: $$\mathcal{E}_{ind}=-\cfrac{\pi{d}^2}{4}\dot{B}_z $$Apply Kirchhoff's second rule to the selected circuit and get:


$$L\dot{I}+IR=-\cfrac{\pi{d}^2}{4}\dot{B}_z $$

\end{solcard}

\begin{solcard}{A2}{Obtain the dependence of the eddy current strength $I(t)$ on time at $B_z(t)=B_\mathrm m\cos\omega t$.}{0.50}

Let's imagine the magnetic field induction in complex form: $$B_z(t)=B_\mathrm me^{i\omega t} $$Then: $$I(i\omega L+R)=-\cfrac{\pi{d}^2B_\mathrm m}{4}\cfrac{de^{i\omega t}}{dt}=\cfrac{\pi{d}^2\omega B_\mathrm me^{i(\omega t-\pi/2)}}{4} $$Since $R+i\omega L=\sqrt{R^2+\omega^2L^2}\cdot e^{i\arctan(\omega L/R)}$, for $I(t)$ find: $$I(t)=I_0\cos(\omega t-\phi)=\cfrac{\pi{d}^2\omega B_\mathrm m}{4\sqrt{R^2+\omega^2L^2}}Re\left(e^{i(\omega t-\pi/2-\arctan(\omega L/R))}\right) $$From:

\par\smallskip\noindent\textbf{Answer:} $$I(t)=\cfrac{\pi{d}^2\omega B_\mathrm m}{4\sqrt{R^2+\omega^2L^2}}\cos\left(\omega t-\cfrac{\pi}{2}-\arctan\cfrac{\omega L}{R}\right) $$
\end{solcard}

\begin{solcard}{A3}{From Gauss's theorem, find the radial component of the field $B_{0\rho}(z)$. Further, in all points, use the result obtained.}{0.30}

Let us apply the Gauss theorem on the flux of magnetic field induction to a cylindrical surface coaxial with the $z$ axis: $$\cfrac{\partial \Phi_\text{base}}{\partial z}+\cfrac{\partial \Phi_\text{side}}{\partial z}=0 $$Since the dependence of $B_z(\rho)$ inside the core can be neglected, the last equation can be write in following form: $$\cfrac{\pi{d}^2}{4}\cfrac{\partial B_{0z}(z)}{\partial z}+\pi d\cdot B_{o\rho}(z)=0 $$From here:

\par\smallskip\noindent\textbf{Answer:} $$B_{0\rho}=-\cfrac{d}{4}\cfrac{\partial B_{0z}(z)}{\partial z} $$
\end{solcard}

\begin{solcard}{A4}{Find the dependence of the axial component of the force $F_z(z,t)$ acting on the ring on the coordinate $z$ and time $t$.}{0.20}

Let us determine the projection of the Ampere force on the axis $z$: $$F_z=-\pi{d}IB_\rho $$Then we have:

\par\smallskip\noindent\textbf{Answer:} $$F_z=\cfrac{\pi^2d^4\omega B_{0z}(z)}{16\sqrt{R^2+\omega^2L^2}}\cfrac{\partial{B_{0z}(z)}}{\partial{z}}\cdot \cos(\omega t)\cos(\omega t-\phi),\quad\phi=\operatorname{arctg}\frac R{L\omega} $$
\end{solcard}

\begin{solcard}{A5}{Find $F_z(z)$.}{0.60}

Let's use the formula for the product of cosines: $$\cos\alpha\cos\beta=\cfrac{\cos(\alpha-\beta)+\cos(\alpha+\beta)}{2}\Rightarrow \cos(\omega t)\cos(\omega t-\phi)=\cfrac{\cos(\phi)+\cos(2\omega t-\phi)}{2} $$Averaging in time leads to the expression: $$\langle\cos(\omega t)\cos(\omega t-\phi)\rangle=\cfrac{\cos\phi}{2}=-\cfrac{\sin\left(\arctan\cfrac{\omega L}{R}\right)}{2}=-\cfrac{\omega L}{\sqrt{R^2+\omega^2L^2}} $$from where we find for the time-averaged force:

\par\smallskip\noindent\textbf{Answer:} $$F_z(z)=-\cfrac{\pi^2d^4B_{0z}(z)}{32}\cfrac{\partial{B_{0z}(z)}}{\partial{z}}\cdot\cfrac{\omega^2L}{R^2+\omega^2L^2} $$
\end{solcard}

\begin{solcard}{B1}{Find an expression for $F_z(z)$ in the geometry under consideration and the potential of this force $U_F(z)$ (assume it is zero at the top of the core).}{0.50}

For the model described in the condition: $$\cfrac{\partial B_{0z}(z)}{\partial z}=-\cfrac{B_0}{H}=const $$Hence obtain expression for force $F_z(z)$:


For the potential taking into account $U_F(H)=0$ we have: $$U_F(z)=-\displaystyle\int\limits_H^zF_z(z)dz=-\cfrac{\pi^2d^4B^2_0}{32H}\cdot\cfrac{\omega^2L}{R^2+\omega^2L^2}\displaystyle\int\limits_H^z\left(1-\cfrac{z}{H}\right)dz=\cfrac{\pi^2d^4B^2_0}{32}\cdot\cfrac{\omega^2L}{R^2+\omega^2L^2}\displaystyle\int\limits_0^{1-z/H}tdt $$from where:

\par\smallskip\noindent\textbf{Answer:} $$F_z(z)=\cfrac{\pi^2d^4B^2_0}{32H}\cdot\cfrac{\omega^2L}{R^2+\omega^2L^2}\left(1-\cfrac{z}{H}\right) $$
\end{solcard}

\begin{solcard}{B2}{At what height $h$ will the ring float in a state of equilibrium if the field is “turned on” slowly? Express your answer in terms of $m$, $g$, $R$, $L$, $\omega$, $d$, $B_0$, $H$.}{0.50}

With a slow change in the current strength in the solenoid, the ring will always be located near the equilibrium position, therefore the height of its levitation is determined by the condition of its equilibrium: $$F_z=mg $$Substitute expression for $F_z(z)$ and account for, that in equilibrium $z=h$, find: $$h=H\left(1-\cfrac{32mgH(R^2+\omega^2L^2)}{\pi^2d^4B^2_0\omega^2L}\right) $$The given answer is applicable if the quantity in the bracket is positive; otherwise $h=0$, therefore finally:

\par\smallskip\noindent\textbf{Answer:} $$h=\begin{cases} H\left(1-\cfrac{32mgH(R^2+\omega^2L^2)}{\pi^2d^4B^2_0\omega^2L}\right)\quad\text{for}\quad \cfrac{32mgH(R^2+\omega^2L^2)}{\pi^2d^4B^2_0\omega^2}<1\\ 0\quad\text{for}\quad \cfrac{32mgH(R^2+\omega^2L^2)}{\pi^2d^4B^2_0\omega^2}>1 \end{cases} $$
\end{solcard}

\begin{solcard}{B3}{To what height $h_\mathrm m$ will the ring jump if the field is “turned on” quickly? Express your answer in terms of $m$, $g$, $R$, $L$, $\omega$, $d$, $B_0$, $H$.}{0.50}

Let us obtain the dependence of the resultant Ampere forces and gravity acting on the ring on its coordinate $z$: $$R_z=F_z-mg=\cfrac{\pi^2d^4B^2_0}{32H}\cdot\cfrac{\omega^2L}{R^2+\omega^2L^2}-mg-\cfrac{\pi^2d^4B^2_0z}{32H^2}\cdot\cfrac{\omega^2L}{R^2+\omega^2L^2}=A-kz $$This linear function, position which correspond value $z_0=h=A/k$, find in previous part, moreover if $A<0$ the ring does not jump, so $h_m=0$. If in position equilibrium potential energy ring in the field, all forces are assumed to equal zero, and the dependence on $z$ has the following form: $$U(z)=\cfrac{k(z-z_0)^2}{2} $$The ring stops at the point at which its potential energy that same, that and in initial position, from which: $$z_\text{stop}=2z_0\Rightarrow h_m=2h=2H\left(1-\cfrac{32mgH(R^2+\omega^2L^2)}{\pi^2d^4B^2_0\omega^2L}\right) $$Obtain answer is applicable onlonly in the case $h_m < H$. For $z > H$ the magnetic field induction $B_z(z) = 0$, and, as consequence, $F_z=0$. Take over zero potential energy in vertex/peak the core, for initial and final potential energy ring obtain: $$U_0=\cfrac{\pi^2d^4B^2_0}{64}\cdot\cfrac{\omega^2L}{R^2+\omega^2L^2}-mgH\qquad U_\text{to}=mg(h_m-H) $$Equate, find: $$h_m=\cfrac{\pi^2d^4B^2_0}{64mg}\cdot\cfrac{\omega^2L}{R^2+\omega^2L^2} $$Final the answer will be written as follows:

\par\smallskip\noindent\textbf{Answer:} $$h_\mathrm m=\begin{cases} 0\quad \text{for}\quad \cfrac{32mgH(R^2+\omega^2L^2)}{\pi^2d^4B^2_0\omega^2L}\geq 1\\ 2H\left(1-\cfrac{32mgH(R^2+\omega^2L^2)}{\pi^2d^4B^2_0\omega^2L}\right)\quad \text{for}\quad \cfrac{1}{2}\leq \cfrac{32mgH(R^2+\omega^2L^2)}{\pi^2d^4B^2_0\omega^2L}\leq 1\\ \cfrac{\pi^2d^4B^2_0}{64mg}\cdot\cfrac{\omega^2L}{R^2+\omega^2L^2}\quad\text{for}\quad \cfrac{32mgH(R^2+\omega^2L^2)}{\pi^2d^4B^2_0\omega^2L}\leq\cfrac{1}{2} \end{cases} $$
\end{solcard}

\begin{solcard}{B4}{Find to what height $h_\mathrm m$ in the general case the ring will jump when the field is turned on quickly, if when it is turned on slowly it levitates at a height $h$. Compute $h_\mathrm m$ if $h=18~\text{cm};8~\text{cm}$. Give your answers to two significant figures.}{0.80}

Let us obtain the connection between $h_\mathrm{m}$ and $h$: at $h=0$ the value $h_\mathrm{m}$ is also equal to zero. In the range $0 < h < H/2$ we obtained $h_\mathrm{m} = 2h$, and for $h \geq H/2$ we obtained $h_\mathrm{m} = H^2 / (2(H-h))$. Since $H = 20~\text{cm}$, we have $h_1 = 18~\text{cm} > H/2$ and $h_2 = 8~\text{cm} < H/2$. Then from the resulting connection between $h_\mathrm{m}$ and $h$ we find: The final dependency $h_\mathrm m(h)$ is as follows: $$h_\mathrm m=\begin{cases} 2h\quad\text{for}\quad 0\geq h\leq \cfrac{H}{2}\\ \cfrac{H^2}{2(H-h)}\quad\text{for}\quad h\geq\cfrac{H}{2} \end{cases} $$

\par\smallskip\noindent\textbf{Answer:} The final dependency $h_\mathrm m(h)$ is as follows: $$h_\mathrm m=\begin{cases} 2h\quad\text{for}\quad 0\geq h\leq \cfrac{H}{2}\\ \cfrac{H^2}{2(H-h)}\quad\text{for}\quad h\geq\cfrac{H}{2} \end{cases} $$
\end{solcard}

\begin{solcard}{C1}{Find the additional average Ampere force $F_z(v_z,z)$ resulting from the motion of the ring.}{0.50}

We express the additional increment of the derivative $\Delta \dot{\Phi}$ associated with the motion of the ring as $$ \Delta \dot{\Phi} = \frac{\pi d^2}{4}\frac{\partial B}{\partial z}\frac{dz}{dt} =  \frac{\pi d^2}{4}\frac{\partial B}{\partial z} \dot{z}. $$This induces an additional current $$ \Delta I = \frac{\pi d^2}{4R}\frac{\partial B}{\partial z} \dot{z}. $$Additional force: $$ \Delta F = B_{\rho} \Delta I \pi d = -\frac{\pi^2 d^4}{16R}\left(\frac{\partial B}{\partial z}\right)^2 \dot{z}. $$Average $\frac{\partial B}{\partial z}$ in time, we obtain an expression for the force.

\par\smallskip\noindent\textbf{Answer:} \[F_z(v_z,z)=-\frac{\pi^2d^4B_0^2v_z}{32 RH^2}\]
\end{solcard}

\begin{solcard}{C2}{Find the damping decrement $\delta$. Express your answer in terms of $R$, $L$, $\omega$, $g$, $H$, $h_\mathrm m$. Calculate $\delta$ to two significant figures.}{1.00}

Since the additional force is directly proportional to the speed, the motion of the ring will be a simple oscillation of a damped oscillator. The damping decrement of such oscillations is equal to:\[\delta=\pi\frac{2\gamma}{\omega_0}.\]In terms of the quantities given in the condition, we can write:\[2\gamma=\frac{\pi^2d^4B_0^2}{32m RH^2},\quad\omega_0^2=\cfrac{\pi^2d^4B^2_0}{32mH^2}\cdot\cfrac{\omega^2L}{R^2+\omega^2L^2}\]Thus, the damping decrement:\[\delta=\frac{\pi^2 d^2B_0}{HR\omega}\sqrt{\frac{R^2+\omega^2L^2}{32mL}}.\]Taking into account the results of the previous part, we express the answer in terms of the required quantities:\[\delta=\pi\frac{R^2+L^2\omega^2}{LR\omega^2}\sqrt\frac{2g}{2H-h_\mathrm m}=0.32\]

\par\smallskip\noindent\textbf{Answer:} Note: As you can see, it actually turned out to be $\delta\sim1$, i.e. In a good way, it is impossible to neglect this EMF. However, this would greatly complicate the calculations without changing the qualitative result.
\end{solcard}

\begin{solcard}{D1}{Find the resistance $R_n$ of a stack of $n$ contact rings. Find the specific force $F_n/nmg$ acting on each of the rings in the stack at $z=0$. At what height $h_n$ does the stack levitate? The rings can be considered very thin, so the inductance of a stack of rings is equal to the inductance of a single ring. Express your answers using $n$, $H$, $R$, $L$, $\omega$.}{0.80}

Since the current strength in the ring is such that it is just beginning to break away from the support: $$\cfrac{32mgH(R^2+\omega^2L^2)}{\pi^2d^4B^2_0\omega^2L}=1 $$Since $R\sim S^{-1}$, and for $n$ rings in contact the total transverse cross-section increase in $n$ times: $$R_n=\cfrac{R}{n} $$Then for force $F_n(z)$ have: $$F_n(z)=\cfrac{\pi^2d^4B^2_0}{32H}\cdot\cfrac{\omega^2L}{(R/n)^2+\omega^2L^2}\left(1-\cfrac{z}{H}\right) $$From here:


The expression for $h_n$ is obtained from the equation: $$\cfrac{F_n}{nmg}\left(1-\cfrac{h_n}{H}\right)=1 $$from:

\par\smallskip\noindent\textbf{Answer:} $$\cfrac{F_n}{nmg}=\cfrac{R^2+\omega^2L^2}{R^2/n+n\omega^2L^2}\left(1-\cfrac{z}{H}\right) $$
\end{solcard}

\begin{solcard}{D2}{At what number of rings $n_\mathrm{opt}$ will the specific force be greatest? Calculate it for the ring considered in the problem.}{0.50}

For any values of $z$, the specific force will be greatest at the lowest value of the denominator. Its minimum possible value is determined by the equality: $$\cfrac{R^2}{n}=n\omega^2L^2\Rightarrow n=\cfrac{R}{\omega L}\approx 3{.}18 $$The integer value of $n$ giving the closest minimum of the denominator equals either $\bigl[R/(\omega L)\bigr]$, or $\bigl[R/(\omega L)\bigr]+1$. Compare these two values: $$\cfrac{1}{3}+3\cdot\cfrac{\omega^2L^2}{R^2}=0{.}629\qquad \cfrac{1}{4}+4\cdot\cfrac{\omega^2L^2}{R^2}=0{.}645 $$Thus:

\par\smallskip\noindent\textbf{Answer:} $$n_{opt}=\biggl[\cfrac{R}{\omega L}\biggr]=3 $$
\end{solcard}

\begin{solcard}{D3}{At what maximum number of rings $n_\mathrm{max}$ will the stack of rings still levitate? Calculate it for the ring considered in the problem.}{0.50}

The condition for possible levitation can be written as follows: $$\cfrac{F_n}{nmg}\geq 1 $$Maximally possible quantity $n$ correspond $z=0$, therefore: $$\cfrac{F_n(0)}{nmg}=\cfrac{R^2+\omega^2L^2}{R^2/n+n\omega^2L^2}\geq 1\Rightarrow n^2-n\left(1+\cfrac{R^2}{\omega^2L^2}\right)+\cfrac{R^2}{\omega^2L^2}\leq 0 $$Maximum value $n$ be determined integer part of the larger root: $$n_{max}=\Biggl[\cfrac{1+R^2/(\omega^2L^2)+\left|1-R^2/(\omega^2L^2)\right|}{2}\Biggr] $$Since $R>\omega L$, we finally have:

\par\smallskip\noindent\textbf{Answer:} $$n_{max}=\biggl[\cfrac{R^2}{\omega^2L^2}\biggr]=10 $$
\end{solcard}

\begin{solcard}{E1}{Find how the average power $R\overline{i^2}$ of the Foucault currents arising in the coil depends on the coordinate $z$. Express your answer in terms of $z$, $H$, $R$, $L$, $\omega$, $g$ and $h$.}{0.70}

For average power we have: $$\langle P\rangle=R\langle I^2\rangle. $$ Expression for $\langle I^2\rangle$ be obtained averaging square cosine: $$\langle I^2\rangle=\cfrac{\pi^2d^4\omega^2B^2_z(z)}{32(R^2+\omega^2L^2)}. $$ Use expression for $h$ from point B2, we get the answer:

\par\smallskip\noindent\textbf{Answer:} $$R\overline{i^2}= \frac{gH^2mR}{L(H-h)}\left(1-\cfrac{z}{H}\right)^2 $$
\end{solcard}

\begin{solcard}{E2}{Find numerically the total Joule loss $Q$ in the ring during its acceleration. Give your answer to two significant figures.}{1.30}

The force acting on the ring depends linearly on $z$: $$ F(z) = F_0 \left(1-\frac{z}{H}\right). $$From the condition $F(h) = mg$ obtain $F_0 = 10 mg$. The equation of harmonic oscillations takes the form: $$ m \ddot{z} =  F_0 \left(1-\frac{z}{H}\right)-mg = -\frac{F_0}{H} z + const. $$Angular frequency of harmonic oscillations: $$ \omega_0 = \sqrt{10\frac{g}{H}} \approx 22.14~\mathrm{s}^{-1}. $$Use initial condition, obtain dependence $z(t)$: $$ z(t) = h(1-\cos(\omega_0 t)). $$Let the ring fly off the core at the moment of time $t_0$: $$z(t_0) = h(1-\cos(\omega_0 t_0)) = H, \quad t_0 = \frac{1}{\omega_0} \arccos \left(\frac{h}{H}-1 \right) \approx 75.5~\text{ms}. $$The power dissipation can be expressed as a function of time by substituting $z(t)$: $$ P(t) = \frac{gH^2mR}{L(H-h)}\left(1-\frac{h}{H}\left(1-\cos(\omega_0 t)\right)\right)^2. $$From here we obtain the expression for $Q$: $$ Q = \frac{gH^2mR}{L(H-h)} \int\limits_0^{t_0} \left(1-\frac{h}{H}\left(1-\cos(\omega_0 t)\right)\right)^2 dt \approx 22.1~\text{J} $$

\par\smallskip\noindent\textbf{Answer:} $$ Q \approx 22~\mathrm{J} $$
\end{solcard}

\begin{solcard}{E3}{Find numerically the mechanical energy $E$ transferred to the ring during its acceleration. Give your answer to two significant figures.}{0.30}

Using the results obtained when solving item B3 for the case $h>H/2$, we obtain: $$E=mgh_m= \cfrac{mgH^2}{2(H-h)}. $$

\par\smallskip\noindent\textbf{Answer:} $$ E=\cfrac{mgH^2}{2(H-h)} \approx 0.29~\mathrm{J} $$
\end{solcard}

\begin{solcard}{E4}{Find numerically the efficiency $\eta$ of the Thomson induction accelerator. Give your answer to two significant figures.}{0.20}

Expression for efficiency: $$ \eta=\frac{E}{E+Q}. $$

\par\smallskip\noindent\textbf{Answer:} $$ \eta \approx 1.3 \text{\%} $$
\end{solcard}

\end{document}